\FloatBarrier
\section{Laboratorio Python: Metriche di Validazione e Selezione di \texorpdfstring{$K$}{K}}
\label{sec:dm-python-lab}

In questa sezione finale vengono presentate le metodologie operative di laboratorio~\cite{tan2018}, evidenziando l'impiego di \texttt{scikit-learn}, \texttt{scipy} e \texttt{pandas} per la preparazione dei dati, l'addestramento dei principali algoritmi di clustering (\textit{K-Means}, \textit{DBSCAN} e \textit{Hierarchical Clustering}), la quantificazione delle metriche interne ed esterne e la stima del numero ottimale di cluster.

\FloatBarrier
\subsection{Preparazione dei Dati e Addestramento dei Modelli di Clustering}
Per eseguire la validazione su un dataset tabellare (ad esempio \texttt{correct\_tips.csv}), si selezionano le feature numeriche (\texttt{total\_bill}, \texttt{tip}, \texttt{size}, \texttt{age}), si applica la normalizzazione Min-Max e si addestrano i modelli di clustering:

\begin{lstlisting}[style=code, language=Python, caption={Normalizzazione Min-Max e addestramento di K-Means, DBSCAN e Clustering Gerarchico.}]
import numpy as np
import pandas as pd
from sklearn.preprocessing import MinMaxScaler
from sklearn.cluster import KMeans, DBSCAN, AgglomerativeClustering
from scipy.spatial.distance import pdist
from scipy.cluster.hierarchy import linkage, dendrogram

df = pd.read_csv('correct_tips.csv', index_col=0)
num_col = ['total_bill', 'tip', 'size', 'age']
df_cat = df.drop(columns=num_col)

scaler = MinMaxScaler()
X = scaler.fit_transform(df[num_col])

# K-Means (K = 3)
kmeans = KMeans(n_clusters=3, n_init=10, max_iter=100, random_state=42)
kmeans.fit(X)

# DBSCAN (eps=0.35, min_samples=8)
dbscan = DBSCAN(eps=0.35, min_samples=8)
dbscan.fit(X)

# Clustering Gerarchico Agglomerativo (Ward) e dendrogramma SciPy
Z = linkage(pdist(X, metric='euclidean'), method='ward')
agg_cluster = AgglomerativeClustering(n_clusters=3, linkage='ward')
agg_cluster.fit(X)
\end{lstlisting}

\FloatBarrier
\subsection{Implementazione Metriche Interne in Scikit-Learn}
Dopo aver addestrato un modello $K$-Means sulla matrice normalizzata $X$, le metriche di compattezza ($\mathrm{SSE}$), separazione (Davies-Bouldin) e concordanza (Silhouette) vengono estratte in modo conciso:

\begin{lstlisting}[style=code, language=Python, caption={Calcolo congiunto di $\mathrm{SSE}$, Davies-Bouldin e Silhouette Score in Python.}]
from sklearn.cluster import KMeans
from sklearn import metrics
from sklearn.metrics import silhouette_score

# Addestramento di K-Means con K fissato
kmeans = KMeans(n_clusters=4, n_init=10, max_iter=300, random_state=42)
kmeans.fit(X)

# 1. Coesione intra-cluster: Inerzia (SSE / WSS)
# Attributo built-in che calcola la somma delle distanze quadratiche dai centroidi
sse = kmeans.inertia_
print(f'SSE (Inertia / WSS): {sse:.3f}')

# 2. Separazione inter-cluster: Davies-Bouldin Score
# Valuta il rapporto di similarita peggiore; ottimale vicino a 0
db_score = metrics.davies_bouldin_score(X, kmeans.labels_)
print(f'Davies-Bouldin Score: {db_score:.3f}')

# 3. Silhouette Score globale: media armonica tra coesione e separazione
# Spettro [-1, 1]; ottimale vicino a +1
sil_score = silhouette_score(X, kmeans.labels_)
print(f'Silhouette Score: {sil_score:.3f}')
\end{lstlisting}

\FloatBarrier
\subsection{Stima del Parametro \texorpdfstring{$K$}{K} Ottimale: Il Metodo del Gomito (Elbow/Knee)}
Come formalizzato nel Capitolo~\ref{chap:dm-kmeans-hierarchical} (\autoref{sec:dm-kmeans-elbow-method}), l'identificazione euristica di $K^*$ si basa sull'individuazione del punto di flesso (\textit{knee point}) in cui il guadagno marginale di coesione interna subisce una marcata decelerazione.

Il listato~\ref{lst:dm-elbow-method} illustra la scansione programmatica dell'intervallo $K \in [2, 40]$ e il tracciamento della curva per individuare visivamente la transizione a gomito:

\begin{lstlisting}[style=code, language=Python, label={lst:dm-elbow-method}, caption={Scansione sistematica e tracciamento della curva del gomito per la selezione di $K^*$.}]
import matplotlib.pyplot as plt

sse_list = []
max_k = 40

# Iterazione sull'intervallo di complessita
for k in range(2, max_k + 1):
    km = KMeans(n_clusters=k, n_init=10, max_iter=100, random_state=42)
    km.fit(X)
    sse_list.append(km.inertia_)

# Identificazione del punto critico di ginocchio
k_star = 10
sse_star = sse_list[k_star - 2]  # Allineamento indice base-zero

# Visualizzazione grafica della curva del gomito
plt.figure(figsize=(9, 5))
plt.plot(range(2, max_k + 1), sse_list, marker='o', color='navy', lw=2)
plt.ylabel('SSE (Inertia)', fontsize=14)
plt.xlabel('Numero di Cluster (K)', fontsize=14)
plt.title('Metodo del Gomito per la Stima di K*', fontsize=16)

# Linee tratteggiate di ancoraggio sul punto ottimale
plt.vlines(x=k_star, ymin=0, ymax=sse_star, colors='crimson', linestyles='dashed', lw=1.5)
plt.hlines(y=sse_star, xmin=0, xmax=k_star, colors='crimson', linestyles='dashed', lw=1.5)
plt.scatter(k_star, sse_star, color='crimson', s=90, zorder=5, label=f'Ginocchio K*={k_star}')

plt.legend(fontsize=12)
plt.grid(True, linestyle=':', alpha=0.6)
plt.show()
\end{lstlisting}

\FloatBarrier
\subsection{Validazione Esterna tramite Tabulazioni Incrociate (Crosstab)}
Per esaminare se i cluster numerici scoperti corrispondano a comportamenti differenziati rispetto a variabili categoriche esterne non utilizzate nel clustering (come il giorno della settimana \texttt{day} o la fascia oraria \texttt{time} nel dataset delle mance \texttt{tips}), si impiega la funzione \texttt{pd.crosstab}:

\begin{lstlisting}[style=code, language=Python, caption={Analisi di contingenza empirica e profilazione esterna dei cluster in Pandas.}]
import pandas as pd

# Creazione della matrice di contingenza tra etichette di cluster e variabile esterna
cross_df = pd.crosstab(kmeans.labels_, df_cat['day'])
print("Matrice di Contingenza Cluster vs Giorno:")
print(cross_df)

# Visualizzazione della distribuzione per cluster
cross_df.plot(kind='bar', stacked=False, figsize=(10, 5), colormap='viridis')
plt.title('Distribuzione della variabile esterna "day" per cluster')
plt.xlabel('Cluster ID')
plt.ylabel('Conteggio Istanza')
plt.grid(axis='y', linestyle=':', alpha=0.7)
plt.show()
\end{lstlisting}

La matrice di contingenza consente all'analista di calcolare immediatamente la purezza e l'entropia condizionata del partizionamento rispetto alle etichette categoriche, completando in tal modo il ciclo di valutazione quantitativa della validità.
